% =============================================================================
%  «ТЕОРИЯ СИСТЕМ» — ствол серии «Архитектура Рассуждения»
%  Блок 6. ВЕРТИКАЛЬ: башня уровней.  Разделы 6.1–6.3 (черновик прозы).
%  Стиль: «выводим и показываем»; репо-якоря — в приложении, не в сносках.
%  Самокомпилируемо: pdflatex blok-6-vertikal.tex (см. память latex-build-windows).
% =============================================================================
\documentclass[12pt,a4paper]{article}
\usepackage{cmap}
\usepackage[T2A]{fontenc}
\usepackage[utf8]{inputenc}
\usepackage[russian]{babel}
\usepackage{paratype}
\usepackage{amsmath,amssymb,amsthm}
\usepackage{listings}
\usepackage[expansion=false]{microtype}
\usepackage[margin=28mm]{geometry}

\newcommand{\rezultat}{\medskip\noindent\emph{Что отсюда следует.}\ }

\title{\textbf{Теория Систем}\\[2mm]\large Блок 6. Вертикаль: башня уровней}
\author{}
\date{}

\begin{document}
\maketitle
\thispagestyle{empty}

\noindent\small\textit{Черновик: разделы 6.1--6.3. Продолжает Блок~5 (целое сверх частей).
Соединение по горизонтали разобрано --- теперь по \emph{вертикали}, к башне уровней. Машинные
якоря --- в Приложении.}\normalsize

\bigskip

\noindent Эмерджентность собирала системы \emph{по горизонтали}. Теперь --- \emph{вертикаль}: уровни
(P1) и их башня. Тезис блока: башня уровней не имеет завершённого «верха», и это P4 --- проявленный
уже на самой структуре, а не в приложениях.

% =============================================================================
\section*{6.1\quad Подъём по уровню как функтор}
% =============================================================================

Подъём по уровню (мы видели его как операцию, §2.7) --- на деле \emph{полный верный функтор}.
Уровень-сдвиг переносит триаду целиком и сохраняет морфизмы (\emph{верность}); а все морфизмы между
поднятыми системами приходят с уровня ниже (\emph{полнота}). Иными словами, P1-иерархия --- это
\emph{башня полных вложений}: каждый уровень целиком и без потерь вкладывается в следующий.

\rezultat
Уровни связаны \textbf{функтором}, а не просто отображением. Иерархия --- не россыпь ступеней, а
\emph{башня полных вложений}, где низшее живёт в высшем без искажения.

% =============================================================================
\section*{6.2\quad Башня без завершённого копредела}
% =============================================================================

Будем \emph{итерировать} подъём: уровень за уровнем. Башня \emph{строго растёт} --- предельного уровня
нет, нет завершённого «верха» как объекта. Это \textbf{P4 на самой иерархии}: вторая нота книги ---
граница --- звучит здесь не в приложении, а на чистой структуре уровней.

И это та же самая граница, что аттрактор-как-role-limit (§4.5) и двойственность вырезать\,/\,слить у
(ко)эквалайзеров (§3.6). Одна граница, услышанная в третий раз.

\rezultat
Вертикаль --- \emph{процесс}, а не завершённая лестница: башня уровней не имеет копредела-объекта.
\textbf{P4 проявлен на самой структуре теории} --- ровно то, что Блок~7 назовёт самоподобием.

% =============================================================================
\section*{6.3\quad Иерархия блокирует парадоксы}
% =============================================================================

У вертикали есть и охранная роль. Уровневая градация (P1, §1.4) не позволяет уровню содержать
\emph{самого себя}: само-референция структурно невозможна. Парадоксы Рассела и Кантора, перенесённые
на типы, здесь не \emph{опровергаются} задним числом --- они просто \emph{не строятся}.

\rezultat
Иерархия --- не только структура, но и \textbf{защита}: при правильной градации парадокс
\emph{невозможен}. Это та же хорошая форма (§2.9), прочитанная вертикально: нарушить уровни --- значит
перестать иметь систему.

\medskip
\noindent\emph{Переход.} Объект, время, целое и вертикаль построены. Следующий блок находит их
\emph{единственную} границу --- и показывает, что она внутри собственного ядра теории.

% =============================================================================
\appendix
\section*{Приложение к Блоку 6 — машинные якоря}
% =============================================================================
\small\raggedright\sloppy
\noindent Всё перечисленное машинно проверено; единственная аксиома ветви --- \texttt{classic} (L3).

\begin{itemize}\itemsep2pt
\item \textbf{§6.1.} подъём \texttt{fs\_lift} как полный верный функтор; перенос триады,
\texttt{lift\_elements\_valid\_up}.\\
\emph{Файл:} \texttt{ERRLevelFunctor.v} (сам \texttt{fs\_lift} --- в \texttt{ERRActualization.v}, §2.7).

\item \textbf{§6.2.} башня без копредела --- \texttt{no\_completed\_tower}, \texttt{tower\_climbs}
(башня строго растёт $=$ P4 на иерархии).\\
\emph{Файл:} \texttt{ERRLevelTower.v}.

\item \textbf{§6.3.} блокировка само-референции --- \texttt{err\_self\_reference\_blocked}; слой
\texttt{Level}, \texttt{level\_lt}, \texttt{level\_lt\_irrefl}.\\
\emph{Файл:} \texttt{TheoryOfSystems\_Core\_ERR.v}.
\end{itemize}
\normalsize

% --- Дальше: Блок 7 (ГРАНИЦА, кульминация) по тезисному плану в Книги/Теория Систем/. ---

\end{document}
